\chapter{Verification and Validation Plan}

\section*{Purpose}
This chapter defines the current verification strategy for the appliance.

\section{Current Verification Goal}
The current phase validates that the external request path, MMIO sequencing, wrapper behavior, imported-core integration seam, and STUB-mode board-bring-up infrastructure are stable enough to support first execution on actual Zynq hardware.

\section{Verification Layers}
\begin{longtable}{p{0.22\textwidth} p{0.68\textwidth}}
\toprule
Layer & Objective \\
\midrule
Documentation review & Confirm consistency across requirements, architecture, interfaces, planning, prompts, and STUB-mode bring-up guidance. \\
Software unit tests & Validate request checking, fake backend behavior, backend selection, probe logic, real-backend file-mapping sanity, and STUB selftest behavior. \\
Wrapper and adapter testbench & Validate reset, digest writes, start-to-busy transition, deterministic completion, signature readback, and all currently supported wrapper-visible engine modes. \\
Imported-core inspection & Confirm that project-owned shim logic is based on the actual imported ML-DSA-OSH sign path rather than assumptions. \\
STUB-mode board bring-up & Prepare safe register-visibility and one-transaction STUB validation on the actual target board. \\
\bottomrule
\end{longtable}

\section{Acceptance Focus for the Current Phase}
The current phase is accepted on interface stability, reproducible STUB-mode bring-up preparation, practical board-facing scripts, and preservation of the local development flow. It is explicitly not accepted on the basis of already-proven execution on real hardware.

\section{Current Test Assets}
\begin{itemize}[leftmargin=*]
  \item Python unit tests for digest validation and fake backend state transitions
  \item Python contract tests for the local client and service path
  \item Python tests for backend selection, probe behavior, file-backed real-backend sanity, and STUB selftest behavior
  \item A SystemVerilog wrapper testbench that covers `STUB`, `CORE\_PLACEHOLDER`, and the documented `MLDSA\_OSH` fallback behavior
  \item Imported-core inspection notes tied to the actual upstream source files
\end{itemize}

\section{Current Limits}
The imported upstream sign path still requires stronger mixed-language validation than the local portable toolchain currently provides. In addition, STUB-mode board validation still depends on actual Zynq hardware, a real AXI base address, and Linux MMIO permissions.

\section{Key Exit Criteria}
\begin{itemize}[leftmargin=*]
  \item Deterministic `STUB` mode remains simulation-backed for regression.
  \item The software stack can perform a real-backend register probe without API churn.
  \item The real selftest path checks the documented STUB signature rule explicitly.
  \item The repository is ready for a follow-on phase that performs actual STUB-mode board execution.
\end{itemize}